\chapter{Conclusion and Next Steps}

The consolidated project defines a technically defensible research platform: distributed geodesic radomes, a hybrid aperture with external VHF Yagi antennas and smaller higher-frequency elements, multiband vector reception, local event processing, precise timing, end-to-end calibration and multistatic fusion. The central design principle is not a universal antenna, but a mechanically integrated family of antennas sharing a common face-level timing, calibration and fusion interface.

As a candidate architectural innovation, the project combines the external Yagis with individually shielded receiving pyramidal cells whose absorber lining remains to be selected. Receive-only operation already establishes a low emission signature relative to active radar: there is no self-generated illumination, and any return caused by the station is only secondary scattering of external waves. The additional hypothesis is to preserve the receive path outside the shell while reducing self-interference, coupling, aperture reradiation and scattering signature. The magnitude of this additional gain must be demonstrated through power balance, shielding effectiveness, OTA patterns and angular RCS; the document does not claim absolute invisibility or operational performance.

The crossed-band Yagi assembly is a proposed mechanical integration solution, not yet a demonstrated innovation. A larger external VHF Yagi and a smaller orthogonally mounted UHF Yagi combine distinct spectral roles in one face-level interface, but they remain independent single-polarization channels and are not combined for polarimetric synthesis. Full polarimetry is reserved for coherent dual-port modules within the same band.

The lowest-risk next step is a three-node VHF/UHF demonstrator with reference and surveillance channels, followed by L/S/C and X/Ku/Ka tiles. HF should proceed as a dedicated program for vector sensors and characteristic platform modes. Territorial expansion must wait for measured coverage, link budgets, geometry, false-alarm performance and metrological stability.
